\documentclass[12pt, a4paper]{article}

% ── Encoding & Língua ──────────────────────────────────────────────
\usepackage[utf8]{inputenc}
\usepackage[T1]{fontenc}
\usepackage[brazil]{babel}

% ── Fontes modernas ────────────────────────────────────────────────
\usepackage{lmodern}
\usepackage{microtype}

% ── Layout ─────────────────────────────────────────────────────────
\usepackage[top=2.5cm, bottom=2.5cm, left=3cm, right=2.5cm]{geometry}
\usepackage{float}
\usepackage{setspace}
\onehalfspacing

% ── Cores ──────────────────────────────────────────────────────────
\usepackage[dvipsnames, svgnames, x11names]{xcolor}
\colorlet{corPrimaria}{SteelBlue4}
\colorlet{corSecundaria}{DarkOliveGreen4}
\colorlet{corDestaque}{Firebrick3}
\colorlet{corAviso}{DarkOrange3}
\colorlet{corFundo}{Gray95}
\colorlet{corBasico}{OliveDrab3}
\colorlet{corIntermediario}{DarkOrange2}
\colorlet{corAvancado}{Firebrick2}

% ── TikZ & PGF ─────────────────────────────────────────────────────
\usepackage{tikz}
\usepackage{pgfplots}
\pgfplotsset{compat=1.18}
\usetikzlibrary{
  shapes.geometric, shapes.multipart,
  arrows.meta, arrows,
  positioning, calc, fit,
  backgrounds, shadows,
  mindmap, trees,
  decorations.pathreplacing, decorations.markings,
  patterns
}

% ── Ambientes especiais ─────────────────────────────────────────────
\usepackage{tcolorbox}
\tcbuselibrary{skins, breakable, theorems}

% ── Matemática & símbolos ──────────────────────────────────────────
\usepackage{amsmath, amssymb, amsthm}
\usepackage{pifont}  % \ding{52} = check

% ── Código-fonte ───────────────────────────────────────────────────
\usepackage{listings}
\lstset{
  basicstyle=\ttfamily\small,
  keywordstyle=\color{corPrimaria}\bfseries,
  commentstyle=\color{gray},
  stringstyle=\color{corSecundaria},
  breaklines=true,
  frame=single,
  rulecolor=\color{corFundo},
  backgroundcolor=\color{corFundo!40}
}

% ── Hiperlinks ─────────────────────────────────────────────────────
\usepackage[colorlinks=true, linkcolor=corPrimaria, urlcolor=corSecundaria]{hyperref}

% ── Cabeçalho / rodapé ─────────────────────────────────────────────
\usepackage{fancyhdr}
\pagestyle{fancy}
\fancyhf{}
\rhead{\small\textcolor{gray}{\nouppercase{\leftmark}}}
\lhead{\small\textcolor{gray}{Deep-Research Imobiliário}}
\rfoot{\small\thepage}
\renewcommand{\headrulewidth}{0.4pt}

% ── Ambientes customizados ─────────────────────────────────────────
\newtcolorbox{definicao}[1][]{
  colback=corPrimaria!8, colframe=corPrimaria,
  fonttitle=\bfseries, title={Ficha técnica}, #1
}
\newtcolorbox{conceito}[2][]{
  colback=corSecundaria!8, colframe=corSecundaria,
  fonttitle=\bfseries, title={#2}, #1
}
\newtcolorbox{exemplo}[1][]{
  colback=Gray95, colframe=gray,
  fonttitle=\bfseries\itshape, title={Exemplo}, #1
}
\newtcolorbox{atencao}[1][]{
  colback=corAviso!15, colframe=corAviso,
  fonttitle=\bfseries, title={Atenção}, #1
}

% ── Comandos de apoio ──────────────────────────────────────────────
% Dificuldade 1–5 em estrelas (compatível com pdflatex via amssymb)
\newcommand{\dif}[1]{%
  \textcolor{corAviso}{%
    \foreach \i in {1,...,5}{\ifnum\i>#1\textcolor{gray!35}{$\bigstar$}\else$\bigstar$\fi}%
  }%
}
\newcommand{\esforco}[1]{\textbf{#1}}
\newcommand{\check}{\textcolor{corSecundaria}{\ding{52}}~}

\title{}
\author{}
\date{}

\begin{document}

% ── Capa ────────────────────────────────────────────────────────────
\begin{titlepage}
  \centering
  \vspace*{1.5cm}
  {\Huge\bfseries\color{corPrimaria} Deep-Research Imobiliário \par}
  \vspace{0.4cm}
  {\Large Decomposição de Projeto \par}
  \vspace{0.3cm}
  {\large Agente de pesquisa de mercado imobiliário\\[2pt]
   \textit{LangGraph $\cdot$ Playwright $\cdot$ Claude $\cdot$ Chroma}\par}
  \vspace{1.2cm}

  % Decoração TikZ — pipeline em miniatura
  \begin{tikzpicture}[node distance=0.35cm, every node/.style={font=\footnotesize\bfseries}]
    \foreach \n [count=\i] in {parse, scrape, normalize, index, analyze, rank, report}{
      \ifnum\i=1
        \node[draw=corPrimaria, fill=corPrimaria!12, rounded corners=3pt,
              minimum height=0.7cm, inner xsep=4pt] (n\i) {\n};
      \else
        \node[draw=corPrimaria, fill=corPrimaria!12, rounded corners=3pt,
              minimum height=0.7cm, inner xsep=4pt, right=of n\j] (n\i) {\n};
        \draw[->, >=Stealth, corPrimaria!70, thick] (n\j) -- (n\i);
      \fi
      \xdef\j{\i}
    }
  \end{tikzpicture}

  \vfill
  {\large Documento gerado em: \today \par}
  \vspace{0.5cm}
  \begin{tikzpicture}
    \fill[corPrimaria!25, rounded corners=8pt] (0,0) rectangle (\linewidth, 2pt);
  \end{tikzpicture}
\end{titlepage}

\tableofcontents
\newpage

% ════════════════════════════════════════════════════════════════════
\section{Visão Geral do Projeto}

O \textbf{Deep-Research Imobiliário} é um agente que recebe uma consulta
(\emph{cidade}, \emph{bairro}, \emph{faixa de preço}, \emph{tipo} e
\emph{operação}), faz \emph{scraping} de portais imobiliários, normaliza e
segmenta os imóveis, opcionalmente os indexa para busca semântica, e produz um
\textbf{relatório de mercado} com estatísticas e ranking de oportunidades —
enriquecido por leitura qualitativa do \textbf{Claude}.

A espinha dorsal é um \textbf{grafo de estado do LangGraph}:
\[
\texttt{parse} \to \texttt{scrape} \to \texttt{normalize}
\to \texttt{index} \to \texttt{analyze} \to \texttt{rank} \to \texttt{report}
\]
com uma \textbf{ramificação condicional} após \texttt{normalize} (curto-circuito
para \texttt{report} quando não há imóveis). Cada nó lê e escreve um estado
compartilhado (\texttt{GraphState}).

\begin{conceito}{Como ler este documento}
O sistema foi decomposto em \textbf{8 partes} implementáveis. O diagrama da
Seção~\ref{sec:dag} mostra as dependências; a tabela da Seção~\ref{sec:tabela}
resume esforço e dificuldade. Cada parte (Seções 4–11) traz objetivo,
pré-requisitos, critérios de aceitação e orientações. Construa seguindo o
\textbf{caminho crítico}: \emph{Fundação} $\to$ \emph{Scraping} $\to$
\emph{Normalização} $\to$ \emph{Inteligência LLM} $\to$ \emph{Relatório} $\to$
\emph{Interfaces}.
\end{conceito}

% ════════════════════════════════════════════════════════════════════
\section{Diagrama de Dependências (DAG)}
\label{sec:dag}

As setas indicam ``depende de''. As partes mais à esquerda/baixo são
fundacionais; a interface de usuário, à direita, é a entrega final.

\begin{figure}[H]
  \centering
  \tikzset{
    parte/.style={
      rectangle, rounded corners=5pt,
      draw=corPrimaria, fill=corPrimaria!12,
      text width=2.9cm, align=center,
      font=\small\bfseries, minimum height=1.05cm, drop shadow
    },
    fund/.style={parte, draw=corSecundaria, fill=corSecundaria!15},
    final/.style={parte, draw=corDestaque, fill=corDestaque!15},
    dep/.style={->, >=Stealth, thick, color=corPrimaria!70}
  }
  \begin{tikzpicture}[node distance=1.5cm and 2.2cm]
    % Fundação
    \node[fund] (p1) {P1\\Modelos \& Estado};

    % Nível 1
    \node[parte, above right=1.2cm and 2.4cm of p1] (p2) {P2\\Grafo \& Config};
    \node[parte, below right=1.2cm and 2.4cm of p1] (p3) {P3\\Scraping};

    % Nível 2
    \node[parte, right=of p3] (p4) {P4\\Normalização \& Segmentação};

    % Nível 3 — dois ramos a partir de P4
    \node[parte, above right=0.4cm and 2.2cm of p4] (p5) {P5\\Busca Semântica};
    \node[parte, below right=0.4cm and 2.2cm of p4] (p6) {P6\\Análise \& Ranking (LLM)};

    % Nível 4
    \node[parte, right=of p6] (p7) {P7\\Relatório};

    % Entrega final
    \node[final, above right=0.6cm and 2.2cm of p7] (p8) {P8\\Interfaces\\CLI \& Streamlit};

    % Arestas (depende de)
    \draw[dep] (p1) -- (p2);
    \draw[dep] (p1) -- (p3);
    \draw[dep] (p1.east) to[out=0,in=180] (p4.west);
    \draw[dep] (p3) -- (p4);
    \draw[dep] (p4) -- (p5);
    \draw[dep] (p4) -- (p6);
    \draw[dep] (p6) -- (p7);
    \draw[dep] (p4.east) to[out=-25,in=200] (p7.west);
    \draw[dep] (p2.east) to[out=0,in=150] (p8.north west);
    \draw[dep] (p7) -- (p8);
    \draw[dep] (p5.east) to[out=0,in=180] (p8.west);

    % Legenda
    \begin{scope}[on background layer]
      \node[draw=gray!40, rounded corners, fill=gray!5, font=\footnotesize,
            align=left, below=1.3cm of p4]
        {\textcolor{corSecundaria}{$\blacksquare$}~Fundação \quad
         \textcolor{corPrimaria}{$\blacksquare$}~Módulo \quad
         \textcolor{corDestaque}{$\blacksquare$}~Entrega final};
    \end{scope}
  \end{tikzpicture}
  \caption{Grafo de dependências das 8 partes. Caminho crítico:
  P1 $\to$ P3 $\to$ P4 $\to$ P6 $\to$ P7 $\to$ P8.}
  \label{fig:dag}
\end{figure}

% ════════════════════════════════════════════════════════════════════
\section{Tabela-Resumo das Partes}
\label{sec:tabela}

\begin{table}[H]
  \centering
  \renewcommand{\arraystretch}{1.4}
  \setlength{\tabcolsep}{6pt}
  \begin{tabular}{c|l|c|c|c}
    \rowcolor{corPrimaria!20}
    \textbf{Parte} & \textbf{Nome} & \textbf{Pré-req.} & \textbf{Esforço} & \textbf{Dificuldade}\\
    \hline
    \rowcolor{corBasico!12}
    P1 & Modelos \& Estado            & ---       & P & \dif{1}\\
    \rowcolor{corBasico!6}
    P2 & Grafo \& Configuração        & P1        & P & \dif{2}\\
    \rowcolor{corAvancado!10}
    P3 & Camada de Scraping           & P1        & G & \dif{5}\\
    \rowcolor{corIntermediario!12}
    P4 & Normalização \& Segmentação  & P1, P3    & M & \dif{3}\\
    \rowcolor{corIntermediario!12}
    P5 & Busca Semântica (Chroma)     & P1, P4    & M & \dif{4}\\
    \rowcolor{corIntermediario!12}
    P6 & Análise \& Ranking (LLM)     & P1, P4    & M & \dif{4}\\
    \rowcolor{corBasico!6}
    P7 & Relatório                    & P4, P6    & P & \dif{2}\\
    \rowcolor{corIntermediario!12}
    P8 & Interfaces (CLI \& Streamlit)& P2, P5, P7& M & \dif{3}\\
  \end{tabular}
  \caption{Esforço: P (pequeno) $\cdot$ M (médio) $\cdot$ G (grande).
  Dificuldade de \dif{1} a \dif{5}.}
  \label{tab:partes}
\end{table}

\begin{atencao}
\textbf{P3 (Scraping) é o gargalo de risco.} Portais usam Cloudflare e mudam o
HTML sem aviso. Comece pelo \texttt{MockScraper} (sintético) para destravar
P4--P8 em paralelo, e trate os scrapers reais como \emph{best-effort}.
\end{atencao}

% ════════════════════════════════════════════════════════════════════
\section{Parte 1 — \textit{Modelos \& Estado (Fundação)}}

\begin{conceito}{Objetivo}
Definir o contrato de dados de todo o sistema: os modelos Pydantic
(\texttt{Query}, \texttt{Listing}, \texttt{Opportunity}, \texttt{MarketAnalysis})
e o \texttt{GraphState} (\texttt{TypedDict}) que trafega entre os nós.
\end{conceito}

\begin{definicao}
\textbf{Arquivos:} \texttt{models.py}, \texttt{state.py}\\
\textbf{Pré-requisitos:} nenhum (raiz do projeto)\\
\textbf{Esforço:} \esforco{Pequeno (P)} \qquad \textbf{Dificuldade:} \dif{1}
\end{definicao}

\subsection*{Critérios de Aceitação}
\begin{itemize}
  \item[] \check \texttt{Query} cobre cidade, bairro, operação (venda/aluguel),
        faixa de preço, tipo (apartamento/casa/ambos), \texttt{max\_paginas},
        \texttt{uf}/\texttt{zona} e \texttt{filtrar\_outliers}.
  \item[] \check \texttt{Listing} guarda portal, URL, preço, área, quartos e
        deriva \texttt{preco\_m2} via \texttt{compute\_preco\_m2()}.
  \item[] \check \texttt{GraphState} declara todas as chaves lidas/escritas
        pelos nós (\texttt{raw\_listings}, \texttt{listings}, \texttt{analyses},
        \texttt{rankings}, \texttt{report}).
\end{itemize}

\subsection*{Orientações de Implementação}
Use \texttt{total=False} no \texttt{TypedDict} (cada nó devolve dicionário
parcial que o LangGraph mescla). Mantenha campos opcionais como
\texttt{Optional[...]}: scrapers reais frequentemente não extraem todos os
atributos.

% ════════════════════════════════════════════════════════════════════
\section{Parte 2 — \textit{Grafo \& Configuração}}

\begin{conceito}{Objetivo}
Montar o grafo LangGraph (nós, arestas, ponto de entrada e a aresta condicional
pós-\texttt{normalize}) e centralizar a configuração via ambiente/\texttt{.env}.
\end{conceito}

\begin{definicao}
\textbf{Arquivos:} \texttt{graph.py}, \texttt{config.py}\\
\textbf{Pré-requisitos:} P1\\
\textbf{Esforço:} \esforco{Pequeno (P)} \qquad \textbf{Dificuldade:} \dif{2}
\end{definicao}

\subsection*{Critérios de Aceitação}
\begin{itemize}
  \item[] \check \texttt{build\_graph()} registra os 7 nós e compila sem erro.
  \item[] \check Aresta condicional após \texttt{normalize}: com imóveis vai a
        \texttt{index}; vazio faz \textbf{curto-circuito} a \texttt{report}.
  \item[] \check \texttt{config.py} expõe \texttt{ANTHROPIC\_API\_KEY},
        \texttt{LLM\_MODEL}, \texttt{USE\_CHROMA}, \texttt{EMBED\_BACKEND} e um
        \texttt{CHROMA\_DIR} \textbf{absoluto} (ancorado na raiz do projeto).
\end{itemize}

\begin{atencao}
O \texttt{LLM\_MODEL} default deve ser o ID \textbf{datado}
(\texttt{claude-haiku-4-5-20251001}); o alias sem data retorna \textbf{404}.
\end{atencao}

% ════════════════════════════════════════════════════════════════════
\section{Parte 3 — \textit{Camada de Scraping}}

\begin{conceito}{Objetivo}
Implementar o \emph{fan-out} assíncrono de scrapers (Playwright) tolerante a
falhas: cada portal roda em paralelo, e um que falhe é ignorado sem derrubar o
pipeline.
\end{conceito}

\begin{definicao}
\textbf{Arquivos:} \texttt{scrapers/base.py}, \texttt{mock.py}, \texttt{olx.py},
\texttt{zapimoveis.py}, \texttt{vivareal.py}, \texttt{quintoandar.py};
\texttt{nodes/scrape.py}\\
\textbf{Pré-requisitos:} P1\\
\textbf{Esforço:} \esforco{Grande (G)} \qquad \textbf{Dificuldade:} \dif{5}
\end{definicao}

\subsection*{Critérios de Aceitação}
\begin{itemize}
  \item[] \check \texttt{BaseScraper.scrape} acumula com \textbf{dedupe por URL}
        e \textbf{para no 1º lote sem novidades}; \texttt{detect\_tipo} classifica
        apartamento/casa.
  \item[] \check \texttt{MockScraper} gera dados sintéticos válidos para destravar
        o pipeline sem rede (modo \texttt{--mock}).
  \item[] \check \texttt{scrape\_node} usa \texttt{asyncio.gather} com
        \texttt{return\_exceptions}; portal que lança é logado e ignorado.
  \item[] \check URLs adaptam \emph{operação} (venda/aluguel) e \emph{tipo};
        OLX/ZAP recebem \texttt{uf}/\texttt{zona}.
\end{itemize}

\subsection*{Orientações de Implementação}
Estado real dos seletores (testado): \textbf{QuintoAndar} via
\texttt{aria-label} do card; \textbf{OLX} via texto do card (paginação real
\texttt{?o=N}); \textbf{ZAP} sujeito a Cloudflare intermitente; \textbf{VivaReal}
bloqueado (403, só scaffold). Use \texttt{scripts/scrape\_debug.py} para
diagnosticar mudanças de HTML.

% ════════════════════════════════════════════════════════════════════
\section{Parte 4 — \textit{Normalização \& Segmentação}}

\begin{conceito}{Objetivo}
Transformar \texttt{raw\_listings} ruidosos em \texttt{listings} confiáveis:
deduplicar, filtrar por preço/quartos/tipo, calcular R\$/m², remover
\emph{outliers} e segmentar por tipo de imóvel.
\end{conceito}

\begin{definicao}
\textbf{Arquivos:} \texttt{nodes/parse.py}, \texttt{nodes/normalize.py},
\texttt{segments.py}\\
\textbf{Pré-requisitos:} P1, P3\\
\textbf{Esforço:} \esforco{Médio (M)} \qquad \textbf{Dificuldade:} \dif{3}
\end{definicao}

\subsection*{Critérios de Aceitação}
\begin{itemize}
  \item[] \check \texttt{parse} normaliza a consulta (corrige min/max invertidos,
        \emph{strip}, tipo).
  \item[] \check \texttt{normalize} reaplica o filtro de preço (o das URLs é
        pouco confiável) e \textbf{descarta imóveis sem preço}.
  \item[] \check \texttt{\_remove\_outliers} corta, por segmento, R\$/m²
        $< 0{,}5\times$ ou $> 2{,}5\times$ a mediana (amostra mínima 8); flag
        \texttt{--manter-outliers} desliga.
  \item[] \check \texttt{segmentar} agrupa por tipo, evitando viés de R\$/m² ao
        misturar casa e apartamento.
\end{itemize}

% ════════════════════════════════════════════════════════════════════
\section{Parte 5 — \textit{Busca Semântica (Chroma)}}

\begin{conceito}{Objetivo}
Persistir um histórico vetorial dos imóveis (Chroma) e permitir consulta
semântica posterior \textbf{sem novo scraping}, com filtros estruturados.
\end{conceito}

\begin{definicao}
\textbf{Arquivos:} \texttt{store.py}, \texttt{nodes/index.py}, \texttt{search.py}\\
\textbf{Pré-requisitos:} P1, P4\\
\textbf{Esforço:} \esforco{Médio (M)} \qquad \textbf{Dificuldade:} \dif{4}
\end{definicao}

\subsection*{Critérios de Aceitação}
\begin{itemize}
  \item[] \check \texttt{index\_node} faz \emph{upsert} por URL na coleção
        \texttt{imoveis} quando \texttt{USE\_CHROMA=true}; senão, \emph{pass-through}.
  \item[] \check \texttt{search.py} consulta o histórico com filtros
        \texttt{--operacao/--tipo/--bairro/--cidade/--min/--max}
        (\textbf{preço vai nas flags}, não no \texttt{--texto}).
  \item[] \check Embeddings locais: \texttt{ollama} (\texttt{nomic-embed-text},
        dim 768) ou \texttt{default} (MiniLM ONNX); mesmo backend para indexar e
        buscar.
\end{itemize}

\begin{atencao}
Trocar o modelo de embedding muda a dimensão: \textbf{apague \texttt{.chroma/} e
reindexe}. O \texttt{.env.example} \emph{não} é carregado pelo dotenv — só
\texttt{.env} tem efeito em runtime.
\end{atencao}

% ════════════════════════════════════════════════════════════════════
\section{Parte 6 — \textit{Análise \& Ranking (LLM)}}

\begin{conceito}{Objetivo}
Calcular estatísticas de mercado por segmento (em Python) e enriquecê-las com
leitura qualitativa e justificativas via Claude — degradando graciosamente
quando não há chave de API.
\end{conceito}

\begin{definicao}
\textbf{Arquivos:} \texttt{llm.py}, \texttt{nodes/analyze.py}, \texttt{nodes/rank.py}\\
\textbf{Pré-requisitos:} P1, P4\\
\textbf{Esforço:} \esforco{Médio (M)} \qquad \textbf{Dificuldade:} \dif{4}
\end{definicao}

\subsection*{Critérios de Aceitação}
\begin{itemize}
  \item[] \check \texttt{analyze} produz uma \texttt{MarketAnalysis} por segmento
        (média, mediana, min/max, R\$/m² médio) + \texttt{observacoes} do LLM.
  \item[] \check \texttt{rank} seleciona Top-N abaixo da média de R\$/m² e gera
        \texttt{justificativa} por oportunidade.
  \item[] \check Sem \texttt{ANTHROPIC\_API\_KEY}, os nós usam
        placeholder/justificativa calculada — \textbf{o pipeline continua}.
  \item[] \check Saída do LLM validada por schema (\texttt{MarketObservations})
        via \texttt{langchain-anthropic}.
\end{itemize}

% ════════════════════════════════════════════════════════════════════
\section{Parte 7 — \textit{Relatório}}

\begin{conceito}{Objetivo}
Montar o relatório Markdown final (puro Python), combinando estatísticas e
ranking por segmento, com rótulos adaptados à operação (Preço/Aluguel,
R\$/m² / R\$/m²/mês).
\end{conceito}

\begin{definicao}
\textbf{Arquivos:} \texttt{nodes/report.py}\\
\textbf{Pré-requisitos:} P4, P6\\
\textbf{Esforço:} \esforco{Pequeno (P)} \qquad \textbf{Dificuldade:} \dif{2}
\end{definicao}

\subsection*{Critérios de Aceitação}
\begin{itemize}
  \item[] \check Renderiza estatísticas e ranking \textbf{por segmento}
        (Top-25 por R\$/m²) em Markdown.
  \item[] \check Lida com o caso vazio (curto-circuito): ``Nenhum imóvel
        encontrado''.
  \item[] \check Mostra \texttt{outliers\_removidos} e adapta rótulos por
        operação (venda/aluguel).
\end{itemize}

% ════════════════════════════════════════════════════════════════════
\section{Parte 8 — \textit{Interfaces (CLI \& Streamlit)}}

\begin{conceito}{Objetivo}
Expor o agente ao usuário: uma CLI parametrizável e uma UI Streamlit, ambas
construindo a \texttt{Query} e executando o grafo.
\end{conceito}

\begin{definicao}
\textbf{Arquivos:} \texttt{cli.py}, \texttt{app.py}\\
\textbf{Pré-requisitos:} P2, P5, P7\\
\textbf{Esforço:} \esforco{Médio (M)} \qquad \textbf{Dificuldade:} \dif{3}
\end{definicao}

\subsection*{Critérios de Aceitação}
\begin{itemize}
  \item[] \check CLI \texttt{python -m deep\_research.cli} aceita
        \texttt{--operacao/--tipo/--uf/--zona/--paginas/--mock} e imprime o
        relatório.
  \item[] \check Console Windows força UTF-8
        (\texttt{sys.stdout.reconfigure} + \texttt{Console(legacy\_windows=False)}).
  \item[] \check \texttt{app.py} (Streamlit) coleta os parâmetros e exibe o
        relatório renderizado.
\end{itemize}

% ════════════════════════════════════════════════════════════════════
\section{Caminho Crítico \& Ordem de Trabalho}

\begin{figure}[H]
  \centering
  \begin{tikzpicture}[node distance=0.45cm, every node/.style={font=\small\bfseries}]
    \foreach \n/\t [count=\i] in {%
      P1/Fundação, P3/Scraping, P4/Normaliz., P6/LLM, P7/Relatório, P8/Interfaces}{
      \ifnum\i=1
        \node[draw=corDestaque, fill=corDestaque!12, rounded corners=3pt,
              minimum height=0.9cm, inner xsep=5pt, align=center] (c\i) {\n\\[-1pt]\scriptsize\t};
      \else
        \node[draw=corDestaque, fill=corDestaque!12, rounded corners=3pt,
              minimum height=0.9cm, inner xsep=5pt, align=center, right=of c\j] (c\i) {\n\\[-1pt]\scriptsize\t};
        \draw[->, >=Stealth, corDestaque!70, thick] (c\j) -- (c\i);
      \fi
      \xdef\j{\i}
    }
  \end{tikzpicture}
  \caption{Caminho crítico do projeto. P2, P5 podem ser desenvolvidas em
  paralelo aos ramos correspondentes.}
  \label{fig:critico}
\end{figure}

\begin{conceito}{Estratégia recomendada}
\begin{enumerate}
  \item \textbf{P1 + P2}: fundação e esqueleto do grafo (rápido, destrava tudo).
  \item \textbf{P3 via Mock primeiro}: implemente o \texttt{MockScraper} antes
        dos scrapers reais — isso libera P4–P7 imediatamente.
  \item \textbf{P4 → P6 → P7}: o miolo de valor (estatísticas, LLM, relatório).
  \item \textbf{P5 e scrapers reais}: paralelizáveis; tratam-se como melhorias
        incrementais de risco isolado.
  \item \textbf{P8}: integra tudo nas interfaces de usuário.
\end{enumerate}
\end{conceito}

\end{document}
